\vfill\pagebreak
\section{A number of values known at run time: slices}
\label{sec:collections:slices}

The length of an array is written in its type, so it must be known while
compiling. A program that keeps the numbers the user types does not know how
many there will be until the user stops. A \textit{slice} is a sequence of
values of the same type whose length is known only when the program runs, and
which can grow.

\lstinputlisting[style=coloredverbatim, caption={Keeping the numbers typed, \textit{examples/collections/numbers.yr}}, label=lst:collections:numbers]{examples/collections/numbers.yr}
\vspace{-10pt}%
\begin{lstlisting}[style=bashVerb, escapechar=@+]
@+\textcolor{teal!80}{alice@dev:\mtilde}@+$ ./numbers
Number (0 to stop): 4
Number (0 to stop): 8
Number (0 to stop): 15
Number (0 to stop): 16
Number (0 to stop): 23
Number (0 to stop): 42
Number (0 to stop): 0
You typed 6 numbers: [4, 8, 15, 16, 23, 42]
\end{lstlisting}

The type of a slice is the type of its elements between brackets, with no
length: \token{numbers} is an \token{[i32]}, on line~4. \token{[]} is the empty
slice, of length 0. Everything else works as for arrays: \token{numbers[i]} is
the element at index \token{i}, counted from 0 and checked when the program
runs, \token{numbers.len} is the length, and \token{for} goes through the
elements, with or without their index.

\subsection{Growing a slice}

The operator \tokennolst{\mtildeascii{}} puts two sequences end to end, and
makes a new slice of their elements, which must have the same type:
\tokennolst{[1, 2] \mtildeascii{} [3]} is \token{[1, 2, 3]}, and neither
operand changes. \tokennolst{numbers \mtildeascii{}= [n]}, on line~10, is
short for \tokennolst{numbers = numbers \mtildeascii{} [n]}: it appends the
elements of the array \token{[n]}, one element here, at the end of
\token{numbers}. The right operand is a sequence, not a single value, so
\tokennolst{numbers \mtildeascii{}= n} is refused.

The operator is not \token{+}, because adding two sequences could as well mean
adding their elements one by one, and \tokennolst{\mtildeascii{}} leaves no
doubt. Since it gives the variable a new value, \tokennolst{\mtildeascii{}=}
needs a variable declared \token{mut}, and \token{dmut} is not needed for it. The
rule of Section~\ref{sec:collections:dmut} is unchanged: \token{numbers[0] = 1;}
would need \token{dmut}.

\subsection{From an array to a slice: \texttt{copy}}
\label{sec:collections:copy}

Square brackets around values always make an array, whatever the type of the
variable that receives it.

\begin{lstlisting}[style=coloredverbatimError, escapechar=@, caption=An array where a slice is expected]
use std::io;

fn main() {
  let odd: [i32] = @\hb{[1, 3, 5, 7]}@; // error
  println(odd);
}
\end{lstlisting}

The compiler says \errmsg{E4095}{incompatible types [i32] and mut [mut i32 ;
  4us]}. The two \token{mut} of the message describe the literal, which the
program may change; the next chapter explains them. What matters is that the
value is an array of four elements, where a slice is expected.

The two are stored differently. The elements of an array are stored in the
variable itself, among the variables of the function that declares it, and
they disappear when that function returns. The elements of a slice are stored
elsewhere, in memory that lasts as long as the slice needs them, and the slice
only records where they are and how many there are. That is why a slice can
grow, and why a function can return one. Making a slice from an array therefore
copies the elements to that other place, and Ymir wants a copy to be written:
\token{copy [1, 3, 5, 7]} is a slice of four elements.

This explanation stays vague on purpose: what \textit{stored in the variable}
and \textit{elsewhere} mean takes knowing how the memory of a program is laid
out. Until Chapter~\ref{chap:layout} gives it, the rule must be accepted as it
is: a slice made from the elements of an array is written with \token{copy}.

\subsection{Slices as parameters}

A parameter of type \token{[i32]} accepts a slice of any length, so one function
serves for all of them. The next listing computes a sum and an average.

\lstinputlisting[style=coloredverbatim, caption={Functions that take a slice, \textit{examples/collections/average.yr}}, label=lst:collections:average]{examples/collections/average.yr}
\vspace{-10pt}%
\begin{lstlisting}[style=bashVerb]
Sum: 98
Average: 14
Weekend average: 22
\end{lstlisting}

An array is not a slice, so it needs a \token{copy} too before it is passed to
\token{average}, on line~21. Without it, the call is refused with
\errmsg{E4226}{the call operator is not defined for \errhole{} and
  \errhole{}}, and the note under it gives the reason,
\errmsg{E4095}{incompatible types [i32] and [i32 ; 2us]}.

Line~12 divides an \token{i32} by a length, which is a \token{usize}. Ymir does
not mix integer types in an operation, and would refuse \token{sum(values) /
  values.len}. \token{cast!i32(x)} converts \token{x} to an \token{i32}: as in
\token{read!i32}, the \token{!i32} names the type wanted. A \token{usize} can hold
numbers too large for an \token{i32}, but the lengths of this book stay far
below two billion, where the conversion would go wrong.

\token{average} has a flaw: the average of an empty slice divides by 0. The
program is then stopped by the system with the message \textit{Floating point
  exception}, although no floating point number is involved.
Section~\ref{sec:collections:options} gives a function a way to say that it has
no answer, and an exercise fixes \token{average} with it.

\subsection{Strings are slices}

Chapter~\ref{chap:types} wrote the type of a string \token{[c8]}. It is indeed a
slice of characters, and everything above applies to it: \token{s.len},
\token{s[0]}, \tokennolst{\mtildeascii{}}, \token{==}, and \token{for c in s}.
\token{read} can read a whole line, when it is asked for a \token{[c8]}. A type
of more than one word goes between braces, so the call is
\token{read!\{[c8]\}}, and returns the line typed, without its line break.

\lstinputlisting[style=coloredverbatim, caption={A string read, joined and gone through, \textit{examples/collections/greeting.yr}}, label=lst:collections:greeting]{examples/collections/greeting.yr}
\vspace{-10pt}%
\begin{lstlisting}[style=bashVerb, escapechar=@+]
@+\textcolor{teal!80}{alice@dev:\mtilde}@+$ ./greeting
Your name: Alice
Hello, Alice! That is 13 characters.
A l i c e
\end{lstlisting}

As Section~\ref{sec:types:encoding} explained, the length of a \token{[c8]} is a
number of bytes, and it is a number of characters only for text in ASCII. A loop
over a \token{[c8]} goes through it byte by byte as well, so a name such as
\textit{Zoé} is not printed as three letters.

\subsection{Part of a slice}

A range between the brackets selects part of a slice. \token{s[a .. b]} is the
slice of the elements from index \token{a} to index \token{b}, \token{b}
excluded, as in the range \token{a .. b}
(Figure~\ref{fig:collections:subslice}). Inside the brackets, \token{\$} stands
for the length of the slice, so \token{s[\$ - 1]} is its last element and
\token{s[2 .. \$]} all of its elements except the first two.

\begin{lstlisting}[style=coloredverbatim]
let primes = copy [2, 3, 5, 7, 11, 13];
println(primes[1 .. 4]); // [3, 5, 7]
println(primes[3 .. $]); // [7, 11, 13]
println(primes[$ - 1]);  // 13
\end{lstlisting}

\begin{figure}[h]
  \centering
  \begin{adjustbox}{max width=\linewidth}
    \begin{tikzpicture}[cell/.style={rectangle, draw=black!60, fill=background!40,
          minimum width=10mm, minimum height=7mm, font=\ttfamily},
        picked/.style={cell, fill=applegreen!25}]
      \foreach \value [count=\i from 0] in {2, 3, 5, 7, 11, 13} {
        \ifnum\i<1
          \node[cell] (c\i) at (\i * 1.0, 0) {\value};
        \else\ifnum\i<4
          \node[picked] (c\i) at (\i * 1.0, 0) {\value};
        \else
          \node[cell] (c\i) at (\i * 1.0, 0) {\value};
        \fi\fi
        \node[cflab, below=1mm of c\i] {\i};
      }
      \node[cflab, left=3mm of c0] {\token{primes}};
      \draw[black!60] ($(c1.north west) + (0.5mm, 1.5mm)$) -- ++(0, 1.5mm)
        coordinate (l) -- ($(c3.north east) + (-0.5mm, 3mm)$) coordinate (r)
        -- ++(0, -1.5mm);
      \node[cflab, above=0.5mm] at ($(l)!0.5!(r)$) {\token{primes[1 .. 4]}};
    \end{tikzpicture}
  \end{adjustbox}
  \caption{\label{fig:collections:subslice} The part \token{primes[1 .. 4]} of a
    slice. It starts at index 1 and stops before index 4, as the range
    \token{1 .. 4} does, so it holds three elements.}
\end{figure}


Taking part of a slice copies nothing: the new slice refers to the same elements
as the first, only fewer of them. The same brackets work on an array. When the
bounds are known while compiling, as in \token{days[0 .. 6]}, the part is an
array of its own, of 6 elements. When they are known only at run time, its
length is unknown, so it can only be a slice, and must be copied,
\token{copy days[0 .. n]}, for the reason given in
Section~\ref{sec:collections:copy}.

\subsection{Expanding part of a slice}
\label{sec:collections:expand_slice}

A slice cannot be expanded whole in general, since its length is known only
when the program runs (Section~\ref{sec:collections:expand_array}). A part of a
slice can, as long as its bounds are known while compiling: \token{s[0 .. 3]}
has 3 elements, whatever the length of \token{s}.

\lstinputlisting[style=coloredverbatim, caption={Expanding part of a slice, \textit{examples/collections/boxes.yr}}, label=lst:collections:boxes]{examples/collections/boxes.yr}
\vspace{-10pt}%
\begin{lstlisting}[style=bashVerb, escapechar=@+]
@+\textcolor{teal!80}{alice@dev:\mtilde}@+$ ./boxes
24
120
Panic in file "boxes.yr", at line 8, in function "boxes::firstVolume" !!!
Aborted (core dumped)
\end{lstlisting}

On line~8, \token{expand sizes[0 .. 3]} passes the first three elements of
\token{sizes} as the three arguments of \token{volume}. \token{sizes} is a
parameter, of any length, but the part has 3 elements, and the compiler knows
it. On line~14, the function receives the slice \token{[4, 5, 6]}, whose first
three elements give 120.

On line~15, the slice has only 2 elements, and the part \token{sizes[0 .. 3]}
does not exist. As for an index, the compiler cannot check it while compiling,
and the program checks it when it runs: it panics.

The bounds must be known while compiling. With bounds that are variables, as in
\token{expand sizes[n .. m]}, the length of the part is known only when the
program runs, and the compiler refuses it with \errmsg{E4257}{unknown length of
  expansion for type mut [i32]}.

\subsection{Slices of slices}
\label{sec:collections:grids}

The elements of a slice can be slices. A grid whose size is known only at run
time, such as a table of \token{rows} rows and \token{cols} columns, can be a
slice of \token{rows} slices of \token{cols} elements each, of type
\token{[[i32]]}. As for an array of arrays, \token{grid[i]} is row \token{i},
and \token{grid[i][j]} its element in column \token{j}.

\begin{lstlisting}[style=coloredverbatim]
let rows = 3;
let cols = 4;
let dmut grid: [[i32]] = [];
for _ in 0 .. rows {
  grid ~= [copy [0 ; cols]];
}
grid[1][2] = 7;
println(grid); // [[0, 0, 0, 0], [0, 0, 7, 0], [0, 0, 0, 0]]
\end{lstlisting}

Each row is a slice of its own, made by a \token{copy}, and its elements are
stored away from the variable (Section~\ref{sec:collections:copy}). Reserving
such a block of memory is called an \textit{allocation}. The grid above takes
one allocation per row, and more for \token{grid} itself, which only records
where each row is and how long it is. The rows are not next to each other in
memory: each lies wherever its block was found
(Figure~\ref{fig:collections:grid_layout}, above).

\cautionbox{\token{copy [copy [0 ; cols] ; rows]} looks like a shorter way to make
  the grid, but it does not make \token{rows} rows. \token{[v ; n]} computes
  \token{v} once, and repeats it \token{n} times: every element of the grid is
  then the same row, and \token{grid[1][2] = 7;} changes the third element of
  every row. The loop above makes a new row at each turn. The next chapter
  explains why, and how \token{dcopy} makes the rows distinct.}

\subsection{A grid in a single slice}
\label{sec:collections:flat_grid}

An allocation takes time, and a grid of a thousand rows takes a thousand of
them. A single slice of \token{rows * cols} elements holds the same numbers
with one allocation: \token{copy [0 ; rows * cols]}. Since the result of
\token{copy} is a slice, whose length is not part of its type, the \token{n} of
\token{[v ; n]} may then be known only at run time.

The program lays the rows out one after the other, and finds the elements
itself. Row 0 takes the indices 0 to \token{cols - 1}, row 1 the next
\token{cols}, and so on. The element at row \token{i} and column \token{j}
comes after \token{i} full rows of \token{cols} elements, and after \token{j}
elements of its own row: it is at index \token{i * cols + j}
(Figure~\ref{fig:collections:grid_layout}, below).

\begin{lstlisting}[style=coloredverbatim]
let rows = 3;
let cols = 4;
let dmut flat = copy [0 ; rows * cols];
flat[1 * cols + 2] = 7;
println(flat); // [0, 0, 0, 0, 0, 0, 7, 0, 0, 0, 0, 0]
\end{lstlisting}

\begin{figure}[h]
  \centering
  \begin{adjustbox}{max width=\linewidth}
    \begin{tikzpicture}[cell/.style={rectangle, draw=black!60, fill=background!40,
          minimum width=10mm, minimum height=7mm, font=\ttfamily},
        ref/.style={cell, minimum width=14mm, font=\footnotesize},
        picked/.style={cell, fill=applegreen!25}]
      \node[ref] (o0) at (0, 0) {row 0};
      \node[ref] (o1) at (0, -0.7) {row 1};
      \node[cflab, above=1mm of o0] {\token{grid}};

      \foreach \value [count=\j from 0] in {1, 2, 3} {
        \node[cell] (a\j) at (2.5 + \j * 1.0, 0.6) {\value};
      }
      \foreach \value [count=\j from 0] in {4, 5, 6} {
        \ifnum\j=2
          \node[picked] (b\j) at (4.5 + \j * 1.0, -1.3) {\value};
        \else
          \node[cell] (b\j) at (4.5 + \j * 1.0, -1.3) {\value};
        \fi
      }
      \draw[cfflow] (o0.east) -- (a0.west);
      \draw[cfflow] (o1.east) -- (b0.west);
      \node[cflab, right=3mm of b2] {\token{grid[1][2]}};

      \foreach \value [count=\i from 0] in {1, 2, 3, 4, 5, 6} {
        \ifnum\i=5
          \node[picked] (f\i) at (\i * 1.0 + 1.0, -3.4) {\value};
        \else
          \node[cell] (f\i) at (\i * 1.0 + 1.0, -3.4) {\value};
        \fi
        \node[cflab, below=1mm of f\i] {\i};
      }
      \node[cflab, left=3mm of f0] {\token{flat}};
      \draw[black!60] ($(f0.north west) + (0.5mm, 1.5mm)$) -- ++(0, 1.5mm)
        coordinate (l0) -- ($(f2.north east) + (-0.5mm, 3mm)$) coordinate (r0)
        -- ++(0, -1.5mm);
      \node[cflab, above=0.5mm] at ($(l0)!0.5!(r0)$) {row 0};
      \draw[black!60] ($(f3.north west) + (0.5mm, 1.5mm)$) -- ++(0, 1.5mm)
        coordinate (l1) -- ($(f5.north east) + (-0.5mm, 3mm)$) coordinate (r1)
        -- ++(0, -1.5mm);
      \node[cflab, above=0.5mm] at ($(l1)!0.5!(r1)$) {row 1};
      \node[cflab, right=3mm of f5] {\token{flat[1 * 3 + 2]}};
    \end{tikzpicture}
  \end{adjustbox}
  \caption{\label{fig:collections:grid_layout} Two ways to keep a grid of 2 rows
    of 3 elements. Above, a slice of slices: \token{grid} records where each row
    is, and each row is a block of memory of its own, wherever it was
    allocated. Below, a single slice: the rows lie one after the other in one
    block, and the element at row 1 and column 2 comes after one full row of 3
    elements and 2 elements of its own row, at index $1 \times 3 + 2 = 5$.}
\end{figure}


The next listing prints a multiplication table of the size the user asks for,
kept in a single slice.

\lstinputlisting[style=coloredverbatim, caption={A grid in a single slice, \textit{examples/collections/table.yr}}, label=lst:collections:table]{examples/collections/table.yr}
\vspace{-10pt}%
\begin{lstlisting}[style=bashVerb, escapechar=@+]
@+\textcolor{teal!80}{alice@dev:\mtilde}@+$ ./table
Rows: 4
Columns: 5
1 2 3 4 5
2 4 6 8 10
3 6 9 12 15
4 8 12 16 20
Row 3, column 4: 12
\end{lstlisting}

The loops of lines~8 to~12 write the element of row \token{i} and column
\token{j} at index \token{i * cols + j}, and the loops of lines~14 to~19 read it
back from there. On line~20, the element of the third row and the fourth
column, counted from 1 by the user, is at row 2 and column 3, counted from 0:
index \token{2 * cols + 3}, 13 here.
